Die Stirnfläche des Balls ist die Kreisfläche $A = \pi d^2/4$. Mit der Luftwiderstandskraft $F_\text{W} = \frac{1}{2}\,c_w\,\rho\,A\,v^2$ und dem Aktionsprinzip $F_\text{W} = m\,a$ (die Gewichtskraft ist im Vergleich dazu klein und wird vernachlässigt) folgt
\begin{align}
 a &= \frac{F_\text{W}}{m} = \frac{c_w\,\rho\,\pi d^2\,v^2}{2\cdot 4\,m}
\end{align}
Mit $c_w = \cwKO$ für eine Kugel (FoTaBe S.~189) und $\rho = \rhLO$ (FoTaBe S.~187):
\begin{align}
 a &= \aTTF \\
   &= \frac{\cwK\cdot\rhL\cdot\pi\cdot\left(\dTT\right)^2\cdot\left(\vTT\right)^2}{8\cdot\mTT} \\
   &= \aTT \approx \result{\aTTP{0}{2}}
\end{align}
Das ist ein Vielfaches der Fallbeschleunigung: Ein schneller Tischtennisball wird sehr stark abgebremst.
